\section{Related Work}
\label{sec:rw}

\paragraph{Environment Synthesis for Single-Agent Training.}
Training an agent with reinforcement learning requires environments that execute, hold state, and verify outcomes automatically, and a recent line of work builds such environments at scale for one agent.
AWM synthesizes code-backed worlds with structured databases and rule-based checkers, thereby delivering a task with executable tools and a verifiable answer~\citep{wang2026agentworldmodelinfinity}.
EnvScaler separates the construction of an environment skeleton from the generation of scenarios and validators, which raises the yield of usable tasks~\citep{song2026envscaler}.
Agent-World pairs environment discovery with a diagnostic loop that keeps tasks near the current ability of the policy~\citep{dong2026agentworld}, and ToolVerse follows tool dependencies to build long-horizon tasks~\citep{zhou2026toolverse}.
Further work adapts existing environments through interface components or learns an environment-design role through self-play~\citep{huang2026envharness,liu2026spade}.
Overall, these systems produce executable and automatically verified tasks for a single agent that observes every tool return itself.
Such a task rarely contains subtasks that several agents can execute simultaneously, and its verifier reports only whether one agent finished.

\paragraph{Multi-Agent Collaboration and Orchestration Training.}
Multi-agent systems differ mainly in how the collaboration structure is decided.
AutoGen organizes collaboration as conversations between agents~\citep{wu2023autogenenablingnextgenllm}, and MetaGPT encodes standardized operating procedures as role-based workflows~\citep{hong2024metagptmetaprogrammingmultiagent}.
In both cases a human authors the structure and the structure stays fixed during execution.
A second group searches the structure automatically, through optimizable agent graphs~\citep{zhuge2024languageagentsoptimizablegraphs}, workflow generation~\citep{zhang2025aflowautomatingagenticworkflow,hu2025automateddesignagenticsystems}, and agentic supernets that select a topology per query~\citep{zhang2025maasagenticsupernet}.
Although search improves the structure, the selected structure still fixes collaboration once execution starts.
A third group updates the model that performs the orchestration.
Representative methods include evolving orchestration for agent collaboration~\citep{dang2025puppeteer}, parallel agent reinforcement learning that credits critical steps~\citep{kimiteam2026kimik25}, and unified agent and tool orchestration with parallel subtask decomposition~\citep{yuan2026paramanager}.
Overall, learned orchestration is established as feasible.
In these studies, the branch structure of the training distribution is inherited from existing benchmarks, thereby remaining neither controlled during construction nor available during reward computation.

\paragraph{Positioning of MAW in Multi-Agent Training.}
In this work, we focus on how the multi-agent environments an orchestrator trains in relate to the signal it trains under.
Within backends that already execute and verify, our goal is to scale orchestration training.
Unlike environment synthesis that targets a single agent, MAW treats branch structure as a construction variable and composes reference task graphs with serial and parallel operators.
Dependencies and independent branches are thereby written into a task on purpose.
Unlike orchestration learning that consumes a fixed task distribution and scores a rollout by terminal success, MAW adopts a more deliberate design by carrying the construction record into training.
Specifically, the schedule of the agent team is scored against the independent subtasks that composition has already verified.
We state what is new here precisely.
Learned orchestration and synthesized environments have both been studied, and synthesizing an environment for the orchestration signal itself has not.
In MAW, one construction supplies the task, its branches, and the ground truth that scores how those branches are used.
